\begin{textsample}{POS Dim 3 – pb – Score 21.00 – pb/pb\_inf\_25.txt}  \label{ex:f3_pos_087}
CAMPO DE EXPERIÊNCIA : “ TRAÇOS, \textbf{SONS}, CORES E FORMAS ” OBJETIVOS DE APRENDIZAGEM E DESENVOLVIMENTO Utilizar \textbf{sons} produzidos por materiais, objetos e instrumentos musicais durante \textbf{brincadeiras} de faz de conta, encenações, criações musicais, festas.
Expressar -se livremente por meio do desenho, \textbf{pintura}, colagem, dobradura, escultura, criando produções bidimensionais e tridimensionais.
Reconhecer as qualidades do \textbf{som} ( intensidade, duração, altura e timbre ), utilizando -as em suas produções sonoras e ao \textbf{ouvir} músicas e \textbf{sons}. ( BNCC, 2017 ) \textbf{SUGESTÕES} DE VIVÊNCIAS QUE FAVORECEM A CONSTRUÇÃO DE EXPERIÊNCIAS PELAS \textbf{CRIANÇAS} \textbf{PEQUENAS} NESTE CAMPO As práticas realizadas na \textbf{Instituição} de Educação \textbf{Infantil} devem possibilitar : - Vivências de situações nas quais utilizem materiais, objetos e instrumentos, no sentido de experimentar \textbf{sons} e ritmos, \textbf{repetindo} -os e recriando -os através de instrumentos existentes e/ou construídos ; - Vivências em que tenham contato com várias possibilidades de uso do \textbf{som}, reconhecendo mudanças quanto à intensidade, altura, tempo de duração, dentre outros aspectos ; - Participação em atividades artísticas da escola, festas e outras manifestações, nas quais instrumentos e \textbf{sons} estejam presentes.
CAMPO DE EXPERIÊNCIA : “ \textbf{ESCUTA}, FALA, PENSAMENTO E IMAGINAÇÃO ” OBJETIVOS DE APRENDIZAGEM E DESENVOLVIMENTO Expressar ideias, \textbf{desejos} e sentimentos sobre suas vivências, por meio da linguagem oral e escrita ( escrita espontânea ), de \textbf{fotos}, desenhos e outras formas de expressão. \textbf{Inventar} \textbf{brincadeiras} cantadas, poemas e canções, criando rimas, alterações e ritmos.
Escolher e folhear \textbf{livros}, procurando orientar -se por temas e ilustrações e tentando identificar palavras conhecidas.
Recontar histórias \textbf{ouvidas} e planejar coletivamente roteiros de vídeos e de encenações, definindo os contextos, os \textbf{personagens}, a estrutura da história.
Recontar histórias \textbf{ouvidas} para produção de reconto escrito, tendo o professor como escriba.
Produzir suas próprias histórias orais e escritas ( escrita espontânea ), em situações com função social significativa.
Levantar hipóteses sobre gêneros textuais veiculados em portadores conhecidos, recorrendo a estratégias de observação gráfica e/ou leitura.
Selecionar \textbf{livros} e textos de gêneros conhecidos para a leitura de um \textbf{adulto} e/ou para sua própria leitura ( partindo de seu repertório sobre esses textos, como a recuperação pela memória, pela leitura das ilustrações etc. ).
Levantar hipóteses em relação à linguagem escrita, realizando \textbf{registros} de palavras e textos, por meio de escrita espontânea. ( BNCC, 2017 ) \textbf{SUGESTÕES} DE VIVÊNCIAS QUE FAVORECEM A CONSTRUÇÃO DE EXPERIÊNCIAS PELAS \textbf{CRIANÇAS} \textbf{PEQUENAS} NESTE CAMPO As práticas realizadas na \textbf{Instituição} de Educação \textbf{Infantil} devem possibilitar : - Convivência com \textbf{crianças} da mesma faixa \textbf{etária}, de outras faixas \textbf{etárias} e com \textbf{adultos}, ampliando seu repertório de comunicação e possibilitando que elas expressem \textbf{desejos}, necessidades, sentimentos e opiniões ; - Vivências com diversos gêneros textuais, tais quais : jornais, literatura \textbf{infantil}, revista em quadrinhos, receitas, e-mails, poesias, rimas, canções etc, favorecendo o conhecimento e o reconhecimento destes gêneros ; - Exploração de \textbf{gestos}, expressões corporais, \textbf{sons}, rimas, por meio de parlendas, poesias, canções, \textbf{livros} de histórias e outros gêneros textuais, ampliando gradativamente sua compreensão da linguagem oral e escrita ; - Participação em \textbf{rodas} de conversas, relatos de experiências, contação e reconto de histórias, elaborando narrativas, imaginando e se expressando de diversas formas ( oral, dança, música, desenho, teatro ) - Acesso a diferentes portadores de texto, apropriando -se da sua diversidade e dos seus usos sociais ; - Acesso a diferentes portadores de textos apropriando -se da sua diversidade e de seus usos sociais, a saber, receitas, rótulos, poemas, cantigas etc criando situações de vivências práticas ; - Experiências de letramento com escrita espontânea. - Vivências de situações que promovam \textbf{escuta} e respeito as opiniões de todos ; - Experiências de \textbf{registro} com letras em situações de \textbf{brincadeiras}, jogos e interações ; - Vivências de situações que possibilitem a expressão da escrita de forma individual e coletiva ; - Vivências de situações de escrita espontânea a partir de situações significativas de produção ;

% matched lemmas: adulto, brincadeira, criança, desejo, escuta, etário, foto, gesto, infantil, instituição, inventar, livro, ouvir, pequeno, personagem, pintura, registro, repetir, roda, som, sugestão
\end{textsample}
